\documentclass[11pt,a4paper]{article}
\usepackage[utf8]{inputenc}
\usepackage[french]{babel}
\usepackage[T1]{fontenc}
\usepackage{lmodern}
\usepackage[margin=1.5cm]{geometry}

\begin{document}

% En-tête
\begin{center}
    {\Huge \textbf{Emma Petit}} \\[0.2cm]
    {\Large \textbf{Ingénieur Data Santé}} \\[0.1cm]
    emma.petit@email.com \ | \ 06 41 88 92 64
\end{center}

\vspace{0.3cm}

% Résumé / Profil
\section*{Profil Professionnel}
\noindent
Ingénieur spécialisé dans le traitement des données de santé. Je conçois des pipelines d'analyse automatisés et robustes pour séquencer l'ADN et identifier des biomarqueurs. Passionné par la recherche biomédicale, j'évolue à l'interface entre la médecine de précision et la data science, tout en respectant strictement les normes de confidentialité des données patients (RGPD, HDS).

\vspace{0.3cm}

% Compétences
\section*{Compétences Techniques}
\noindent
\textbf{NGS} $\bullet$ \textbf{Statistiques} $\bullet$ \textbf{Analyse Génomique} $\bullet$ \textbf{Nextflow} $\bullet$ \textbf{BioPython} $\bullet$ \textbf{Bases de données biologiques} $\bullet$ \textbf{Python} $\bullet$ \textbf{R}

\vspace{0.3cm}

% Expériences
\section*{Expériences Professionnelles}
\noindent \textbf{Ingénieur Data Santé / Confirmé} --- \textit{Inserm} \hfill \textbf{11/2020 à Aujourd'hui} \\
Modélisation structurelle 3D de protéines et simulation de dynamique moléculaire, permettant de valider l'affinité de nouvelles molécules thérapeutiques. Développement de pipelines bio-informatiques reproductibles avec Nextflow et Docker, réduisant le temps d'analyse des génomes de 3 jours à 12 heures. Gestion et optimisation des requêtes sur un cluster de calcul haute performance (HPC) sous Slurm pour le traitement de pétaoctets de données brutes. Automatisation du nettoyage et du contrôle qualité (QC) des données issues des séquenceurs Illumina, réduisant les erreurs d'alignement de 15\%. \\[0.3cm]
\noindent \textbf{Ingénieur Data Santé} --- \textit{Institut Pasteur} \hfill \textbf{08/2018 à 09/2020} \\
Automatisation des rapports de diagnostic génétique à l'aide de scripts RMarkdown, garantissant une traçabilité totale pour les audits de santé. Analyse de larges jeux de données de séquençage (NGS) à l'aide de scripts Python et R, aboutissant à l'identification de nouveaux biomarqueurs tumoraux. Conception d'une API RESTful sécurisée pour interroger les bases de données génomiques depuis les applications cliniques partenaires. Conception d'une base de données relationnelle sécurisée (HDS) centralisant les données phénotypiques et génotypiques d'une cohorte de 10 000 patients. \\[0.3cm]
\noindent \textbf{Stage — Assistant Ingénieur Data Santé} --- \textit{Inserm} \hfill \textbf{05/2016 à 03/2018} \\
Rédaction de rapports d'analyse statistique détaillés avec RMarkdown pour orienter les décisions lors des essais cliniques de phase 2 du laboratoire. Croisement et intégration de multiples bases de données biologiques publiques (NCBI, Ensembl, ClinVar) pour annoter des variants génétiques rares. Création d'applications web interactives avec R Shiny permettant aux biologistes d'explorer visuellement les résultats d'expression génique (RNA-Seq). \\[0.3cm]


\vspace{0.3cm}

% Formations
\section*{Formation \& Diplômes}
\noindent \textbf{Licence Professionnelle en Data Science pour la Santé} \hfill \textbf{2017 - 2018} \\
\textit{Université d'Évry} \\[0.3cm]
\noindent \textbf{BTS en Data Science pour la Santé} \hfill \textbf{2015 - 2017} \\
\textit{IUT Génie Biologique} \\[0.3cm]


\end{document}